\scanpage{6}
% archive foliation: 5
Prouvons $\Longleftarrow$, i.e. \deleted{Soit $x\in\overline C(j)$}, $R(x)\ne\varnothing$, $\forall y\in R(x)$, on a $y\in G(j)$, prouvons $x\in G(j)$.

Stratégie $\Sigma$ avec :
\[
z\in C(j),\quad\Sigma(z)=\begin{cases}R(z)\text{ si }z\notin G(j)\\ R(z)\cap G(j)\text{ si }z\in G(j)\end{cases}
\]
 a) \uline{\unsure{Bien une} stratégie}, \deleted{i.e. $R(z)\ne\varnothing\Longrightarrow\Sigma(z)\ne\varnothing$.}\inserted{par $1^\circ$. et $2^\circ)$, $\Longrightarrow$}

\deleted{En effet, il n'y a de question que si $z\in G(j)$ et alors $R(z)\cap G(j)\ne\varnothing$ par a)}

 b) Elle est $x$-gagnante. En effet soit $(x_0=x,x_1=y,x_2,\ldots,x_n,\ldots)$ une partie partant de $x$, comp. avec $\Sigma$, admettons qu'elle est finie.
\margin{On suppose qu'il n'y a \unsure{pas de|de} partie infinie}

Prouvons au préalable \unsure{en prouvant} que $x_i\in G(j)\quad\forall i\geq1$ \deleted{\unsure{les distinguons}} \inserted{OK} pour $i=1$ par hyp.

Supposons que \deleted{\unsure{c'est pour}} l'on ait $x_i\in G(j)$, prouvons $x_{i+1}\in G(j)$. En effet, si $x_i\in\overline C(j)$, cela provient de la partie déjà prouvée de $2^\circ$.
Si $x_i\in C(j)$, cela provient de la définition de $\Sigma$, car $\Sigma(x_i)=R(x_i)\cap G(j)$ donc $x_{i+1}\in\Sigma(x_i)$ \deleted{\unsure{donc x}} est $\in G(j)$. O.K.

\uline{Corollaire} \deleted{\unsure{Pour tte une stratégie}}\deleted{Il existe une stratégie $\Sigma$ telle}\deleted{\unsure{soit} $\Gamma\subset G(j)$.}
\inserted{Soit $j\in J$, $\Gamma\subset G(j)$ stable (dans $G(j)$) par succession.}
a) Il existe une \deleted{\unsure{que}} stratégie $\Sigma$ gagnante pour tous les $x\in\Gamma$.

 b) Parmi celles-ci il y en a une plus grande
